\section{Motivation for a Modeller's Toolkit}

The analysis of seismic hazard and risk represents a complex synthesis of many elements of earthquake science and engineering. The knowledge of earthquake behaviour is effectively distilled into a framework in which effective models of the earthquake process can be integrated into a probabilistic calculation, the output of which can provide a basis for effective decision making in the mitigation process. The Global Earthquake Model (\gls{acr:gem}) initiative has many objectives toward improving, on a global scale, the state of knowledge and state of the science that fundamentally underpins the seismic hazard and risk analysis process. It is in addressing these objectives that \gls{acr:gem} has undertaken of development of OpenQuake, and open-source integrated platform for the analysis of seismic hazard and risk. Whilst the OpenQuake platform itself contains within it a powerful set of tools for implementing seismic hazard and risk calculations, it requires as input a fully defined seismic input model for seismogenic sources and ground motion prediction equations (the details of the OpenQuake itself are described \cite{crowley2011}, and we refer the reader to that publication for more information on some of the concepts described within this document). In practice the seismic hazard analysis procedure begins sooner, at the the model development stage. 

It is at this point that it is helpful to clarify what is meant by the term "`hazard model development process". In earthquake science, there are many different sources of information about the nature of earthquake behaviour, from observed seismicity (either recorded by seismic instrumentation [instrumental seismicity] or by historical observation and accounts of past events [historical seismicity]), seismoteconic parameters (e.g. earthquake focal mechanisms), active fault and geodetic observation. The "`seismogenic source model"' itself, however, must define each seismic source by a geometry (point, polygon or plane), a model describing the probability of the source generating an earthquake with a given magnitude (recurrence model), a characterisation of the type of faulting expected from the source (the source mechanism) and additional constraints on the manner in which certain elements of the earthquake source may vary within the model. The attributes of a seismogenic source model represent the output of a procedure of analysis and interpretation. 

The methodology for calculation of probabilistic seismic hazard has become generally well-established since its inception by \cite{cornell1968} and \cite{mcguire1976}, and whilst there have been additional improvements to many aspects of the methodology, the fundamentals processes remain. For the development of the PSHA input model, a "core" methodology has emerged from the state of practice yet in many instances fundamental modelling decisions may be applied without recognition of the impact that those decisions may have on the hazard model. Furthermore, at each step of the model building process it is often that case that different possible algorithms or methodologies could be considered, not all of which could be used interchangeably. The picture becomes complicated yet further when taking into consideration the different possible ways in which data and concepts may be defined, particularly with respect to the issues of uncertainty and data quality. 


In the development of a global model for seismic hazard and risk it is necessary to understand the different options that are available and the possible impact that they have on the model development process. It is in this context that the "\gls{acr:gem} Seismic Hazard Modeller's toolkit"' is developed. Whilst at present it remains in an early stage of development, the objective of the Modeller's Toolkit is to provide a computational tool that gives the user a means of implementing many of the current methodologies applied in hazard model development. It does so within a single computational framework (discussed in due course), allowing the user the ability to utilise the tools directly and where possible, chain together multiple steps to provide an end-to-end model development process. Furthermore, in being developed within the scope of \gls{acr:gem}, the toolkit aims to make the implementation of the tools transparent and reliable by virtue of the test-driven open-source development process.

This document serves a dual purpose as to provide the user with both an explanation of the methodologies currently implemented in the toolkit (including discussion of future tools that are currently in planning and development) and as a manual to allow the reader to install, run and interpret the Modeller's Toolkit tools. 